%=============================================================================!
% 10.S  WHAT WE NOW HAVE                                                      !
%=============================================================================!
\unnumbered{What we now have}

A sample small enough to simulate is mostly surface, and repeating a cell
of atoms in every direction removes the surface. The cell is the matrix
$\vb{h}$ of its lattice vectors; fractional coordinates $\vb{s} =
\vb{h}^{-1}\vb{r}$ wrap positions into it and, rounded, give the nearest
copy of a separation, in any cell, for every separation shorter than half
the smallest perpendicular width $w_{\min}$, the least of widths such as
$V/\lvert\vb{b}\times\vb{c}\rvert$, which bounds the cutoff.

Charges in a periodic cell need the Ewald sum: $1/r$
split by Gaussian integrals into a short-ranged part summed over near
copies and a smooth part summed as a Fourier series over the reciprocal
vectors, less the self term, with the $\vb{G} = \vb{0}$ term dropped for
a neutral cell; it reproduces the Madelung constant of rock salt for any
splitting and any equivalent cell.

Cell lists find the pairs within the
cutoff at a cost that grows as $N$, and Verlet lists with a skin reuse
them until an atom has moved half the skin. A run starts from a crystal
at its model's own spacing, with random velocities stripped of drift, and
the loop is velocity Verlet over a model that returns energy and forces.
The code agrees with ASE to the rounding of the arithmetic once both use
the same physical constants, conserves energy and momentum in a periodic
box, and passes the $\dt^2$ test with Ewald forces. A profiler found the
time in the neighbour list; the force loop, compiled by Numba or written
in Fortran, runs hundreds of times faster than the Python loop that does
each pair's arithmetic with NumPy, and agrees with it to rounding.

Chapter~11 holds the fastest vibrations fixed with constraints, so that
the step can be longer, and splits the forces by how fast they change.
